\section{怎样划分一段话的层次}

我们如果已经弄清怎样顺着思路把紧扣着的一句话又一句话组织成文理通顺的一段话，那么，怎样把一段话剖析成为几个相对独立的层次就比较好办了。

一般说来，一段话有一个集中的意思（段意），而这段意义多不是由一句话表达出来，而是由许多句话，从不同的方面表达出来的。这种情况在比较长的语段中显得尤为突出。

所谓划分一段话的层次，就是把从不同的方面说明段意而彼此之间的意思相对独立的句子或句群（一组意义上联系紧密的句子）划分开来。因此，划分一段话的层次，必须首先统观全段文字，弄清这段话集中表达什么意思；其次再从一句话又一句话分析起。如果这一句话就能表达这段话的一方面意义，而且在意思上与别的句子相对独立，那么这一句
话就是一层意思；如果这一句话不能表达这一方面意思，或者即使能表达这一方面意思而与另外的句子意义紧密相连，不能相对独立，那就要再分析下去，直至这一句话与别的一句或几句意义紧密相连的句子组成一个完整的意思而又能相对独立时为止，我们就把这个句群看作另一层意思。这样一层一层地分析下去，就能把一段话清楚而准确地划分为若干个层次；在划分层次的同时，我们也就能将层意归纳出来。例如高考试卷中有过这样一道题目，它要求考生将下面这段文字划分成四个层次，在每两个层次的交接处划一条竖线，把它们分开，规定一共划三条竖线，少划了或多划了都算错。这段话的文字是这样的（试卷上没加标点）：

永定河上的芦沟桥，修建于公元一一八九到一一九二年间。桥长二百六十五米，由十一个半圆形的石拱组成，每个石拱长度不一，自十六米到二十一点六米。桥宽约八米，路面平坦，几乎与河面平行。每两个石拱之间有石砌桥墩，把十一个石拱联成一个整体。由于各拱相联，所以这种桥叫做联拱石桥。永定河发水时，来势很猛，以前两岸河堤常被冲毁，但是这座桥从没出过事，足见它的坚固。桥面用石板铺砌，两旁有石栏石柱。每个柱头上都雕刻着不同姿态的狮子。这些石刻狮子，有的母子相抱，有的交头接耳，有的象倾听水声，千态万状，惟妙惟肖。
\begin{flushright}
——茅以升《中国石拱桥》
\end{flushright}

我们先统观全段的集中意思在于介绍芦沟桥的概况。尽管全段原文没加标点，但是哪儿可加句号是不难确定的。第一句话说明芦沟桥修建的时间，它既从一个方面说明段意，又与后面说明桥长的句子在意义上相对独立，因此我们应在“一九二年间”之后划一条竖线，表示这句话就是这段话的第一层次。下面几句话，分别说明桥长、桥宽和石拱之间的桥墩的情况，而“由于各拱相联，所以这种桥叫做联拱石桥”这句话是对上面几句话的总结。由于它们相对独立地介绍了这种联拱石桥的具体建筑情况，因此可以在“叫做联拱石桥”的后面划一条竖线，表示这是第二层次。下面一句话是用具体事实说明桥的坚固性，它既从一个方面表达了段意，又与前后的句意相对独立，因此可在“足见它的坚固”的后面再划一条竖线，表示这是第三个层次。最后三句，具体说明桥在建筑上的高度艺术价值，当然是第四层了。

应该注意：一段话中的层次体现了作者对一段话的观点、材料所作的布局和安排，所以在划分一段话的层次时，必须分析这段话的思想脉络和它的结构方式。因此，在划分段落的层次时，一方面不能太琐碎，以免割裂思路和结构方式，如不必把上面那段文字中的桥长、桥宽和桥墩的情况分别划层，另一方面也不能太笼统，以免思路不清，结构不明，如不能把上面那段文字中的第二、三层作为一层看。总之，我们要把能否比较集中地表达某一方面段意而又与其他各层意思相对独立作为是否准确地划分一段话层次的标准。

当然，上面说的是划分语段层次的一般方法。由于文章的文体不同，作者在构段时的思路和结构方式也会因文而异。我们还需进一步掌握各类文体的构段思路和方式，才能做到准确地划分语段的层次。

先说说怎样划分议论文语段的层次。

1.找出段中主句。一般说来，议论文的段中主句就是论点。它往往在一段话的开头出现，用以总领全段，提挈下文。有时它放在一段话的结尾部分，用以总结全段，收束上文。少数文章的段中主句在一段话的中间部分，有总上束下的作用。总之，这个段中主句是一段话的纲，表明了一段话的中心意思。我们应着力找出它。

2.弄清结构方式。议论文语段的结构方式，如在《怎样组织好一段话》的第三点中所说的那样，也大致有三种：一种是“总分式”（或“分总式”、“先总后分再总”式），一种是“并列式”，一种是“推进式”。若是“总分式”的话，则要分析哪几是总述，哪几是分述，又是采用什么方式分述的；若是“并列式”，则要分析其从哪几方面论述；若是“推进式”，则要分析其如何逐层论述。有时，一段话中采用了几种结构方式，比如先总述，再分述，分述时可以先采用“并列式”，然后又用“推进式”，最后再总述，象这样一类复杂的结构方式，我们更应细心分析。由于议论文语段的层次之间逻辑性强，所以并不难弄清。例如：

（1）白求恩同志毫不利己专门利人的精神，表现在他对工作的极端的负责任，对同志对人民的极端的热忱。每个共产党员都要学习他。不少的人对工作不负责任，拈轻怕重，把重担子推给人家，自己挑轻的。一事当前，先替自己打算，然后再替别人打算。出了一点力就觉得了不起，喜欢自吹，生怕人家不知道。对同志对人民不是满腔热忱，而是冷冷清清，漠不关心，麻木不仁。这种人其实不是共产党员，至少不能算一个纯粹的共产党员。从前线回来的人说到白求恩，没有一个不佩服，没有一个不为他的精神所感动。晋察冀边区的军民，凡亲身受过白求恩医生的治疗和亲眼看过白求恩医生的工作的，无不为之感动。每一个共产党员，一定要学习白求恩同志的这种真正共产主义者的精神。
\begin{flushright}
——毛泽东《纪念白求恩》
\end{flushright}

本段层次划分如下：“········每个共产党员都要学习他。$\vert$不少的人········不能算一个纯粹的共产党员。$\vert$从前线回来的人········无不为之感动。$\vert$每一个共产党员········共产主义的精神。$\vert$”第一层是全段的论点，即段中主句；第二层和第三层是采用反面与正面对比的“并列式”结构方式进行论证的，第四层是结论。全段是“先总后分再总式”的结构。

（2）在我国历史上，许多民族英雄、人民英雄、大发明家、科学家，他们都是一些有伟大理想的人。当强敌压境、国家民族危在旦夕的时候，民族英雄的理想，就是要把敌人赶走，使自己的民族生存和发展下去。当统治者昏庸腐朽，横征暴敛，使得人民无法生活下去的时候，人民英雄就揭竿而起，把反抗强权，救民于水火之中作为自己的理想。当时人民的劳动强度很大，生活很苦，劳动生产率很低，发明家、科学家们的理想，就是要以他们的创造、发明，去改善人民的劳动条件，提高劳动生产率和改善人民的生活。归根结底，这些人对促进社会的进步，对社会生产力的发展，是有所贡献的，虽然他们的贡献还不免要受着历史条件的限制。
\begin{flushright}
——陶铸《崇高的理想》
\end{flushright}

本段层次划分如下：“······有伟大理想的人。$\vert$ 当强敌压境······生存和发展下去。$\vert$ 当统治者昏庸腐朽······作为自己的理想。$\vert$ 当时人民······改善人民的生活。$\vert$ 归根结底······受着历史条件的限制。$\vert$  ”第一层是全段的论点，即段中主句，指出我国历史上的伟大人物都是有伟大理想的人；第二、三、四层采用平行叙述的“并列式”结构，分述民族英雄、人民英雄和发明家科学家的理想；第五层采用“推进式”结构，总结他们的贡献和局限。

总之，议论文重在剖析事理，因此，几层意思先后顺序的安排往往以客观事理的逻辑联系为依据。我们只要找准了论点，弄清了合乎逻辑的结构方式，划分议论文语段的层次就不难达到“准确”了。

现在再说说怎样划分记叙文语段的层次。

1.对于“先总后叙”的语段，先确定那总写一笔的“中心句”，再分析具体叙述或描写的部分。“中心句”往往在首句，用以统领全段。具体叙述或描写的部分，则以时间的推移、空间的转换、人物的变化或者认识的发展等为依据来划分层次。例如：

再往里走，天山越来越显得优美。$\vert$ 在那白皑皑的群峰的雪线以下，是蜿蜒无尽的翠绿的原始森林，密密的塔松象无数撑天的巨伞，重重叠叠的枝丫间，只漏下斑斑点点细碎的日影。骑马穿行林中，只听见马蹄溅起在岩石上漫流的水的声音，更增添了密林的幽静。在这林海深处，连鸟雀也少飞来，只偶然能听到远处的几声鸟鸣。当你下马坐在一块岩石上吸烟休息时，虽然林外是阳光灿烂，而在这遮住了天日的密林中却闪着烟头的红火光。从偶然发现的一棵两棵烧焦枯树看来，这里也许来过辛勤的猎人，在午夜生火宿过营，烤过猎获的野味。这天山上有的是成群的野羊、草鹿、野牛和野骆驼。
\begin{flushright}
——碧野《天山景物记》
\end{flushright}

在这段话中，第一句“再往里走，天山越来越显得优美”是中心句，下面分别从“骑马穿行林中”和“下马”休息所见这两个方面描写“越来越显得优美”的景色。明确了中心句与下文的关系，进而又明确了下文两个方面的关系，那么，怎样准确地划分这个语段的层次也就不难了。

2.对于没有中心句，只是叙述事情经过的语段，则可按照情节发展的时间顺序、空间变换、人物变化、作者思路或材料性质等来划分层次。例如：

她不是鲁镇人。 $\vert$ 有一年的冬初，四叔家里要换女工，做中人的卫老婆子带她进来了，头上扎着白头绳，乌裙，蓝夹袄，月白背心，年纪大约二十六、七，脸色青黄，但两颊却还是红的。卫老婆子叫他祥林嫂，说是自己母家的邻舍，死了当家人，所以出来做工了。 $\vert$ 四叔皱了皱眉，四婶已经知道了他的意思，是在讨厌她是一个寡妇。但看她模样还周正，手脚都壮大，又只是顺着眼，不开一句口，很象一个安分耐劳的人，便不管四叔的皱眉，将她收下了。$\vert$ 试工期内，她整天的做，似乎闲着就无聊，又有力，简直抵得过一个男子，所以第三天就定局，每月工钱五百文。
\begin{flushright}
——鲁迅《祝福》
\end{flushright}

这段话的第一句，交代祥林嫂不是鲁镇人。那么她是怎么来的呢？下面接着叙述她被卫老婆子带来鲁镇介绍做工的情形。她是否能被接受做工呢？于是就通过对四叔和四婶的心理描写，交代她被留下试工。试工结果怎样呢？写她终于被正式留下做工了。这里有时间的转移、空间的变换、人物的变化，显示出作者明晰的思路。我们把这些综合起来看，就可以准确地划分出这个语段的层次。

至于如何划分说明文语段的层次，比起如何划分议论文或记叙文语段的层次来，是要简单得多。一般说来，只要按照说明事物的顺序就能划分说明文语段的层次，其结构方式与议论文的三种主要结构方式大致相同。例如开头所引的关于芦沟桥的那段话，其说明顺序就是：
①\ 建桥时间；
②\ 桥的建筑概况；
③\ 桥的坚固性；
④\ 桥的艺术价值。那么，要划分这段话的层次，就该按照上述的说明顺序才行。再如：

新建的北京车站，这座宏伟的建筑物耸立在东单和建国门之间，显示着人民首都的大门所应有的色彩和光辉。$\vert$ 车站大楼巍峨矗立，正面分中部和两翼三个部分。$\vert$ 中间顶端是一个新颖的大扁壳屋顶，屋顶的两旁对称地矗立着两座具有浓厚民族风格的钟楼。$\vert$ 钟楼的屋顶是用金黄色琉璃瓦盖的，表现了新结构同民族传统风格巧妙结合的新颖的建筑艺术。$\vert$ 旅客最关心时间，钟楼下的两座四面大电钟，每面直径长达四米。钟面是墨玉大理石铺贴的，长短针和字标是用乳白色玻璃和金属框子做的。针和字标里面装着电灯，距离三、四公里就能看到，同时还能听到大钟报时报刻的声音。
\begin{flushright}
——陈登攀《首都的大门——北京车站》
\end{flushright}

这段说明性文字采用了由大到小、层层深入的“推进式”结构方式。第一层说明北京车站的位置和价值；第二层说明车站主体大楼的立面结构；第三层说明中间建筑的特点；第四层说明钟楼建筑的特色；第五层说明钟楼上两座四面大电钟的优异构造。我们弄清了这段语的结构方式和说明次序，它的层次也就一目了然了。

这儿所说的划分一段话层次的全部方法，对于文言文来说，也是完全适用的。
\newpage